\section{Техническое задание}

\subsection{Основание для разработки}

Основанием для разработки является задание на выпускную квалификационную работу бакалавра. Тема работы -- <<Разработка программно-информационной системы для управления проектами>>.

Система разрабатывается в рамках преддипломной практики. Она автоматизирует ведение проектов, декомпозицию работ на задачи, управление составом команды и строит диаграмму Ганта как основной способ отображения плана. Требования сформированы по итогам анализа предметной области.

\subsection{Цель и назначение разработки}

В рамках работы разрабатывается веб-ориентированная информационная система управления проектами. В ее составе -- реестр проектов, декомпозиция работ на задачи, управление зависимостями и совместная работа участников с разграниченными правами доступа.

Система предоставляет руководителю и участникам единое рабочее пространство: планирование задач, контроль сроков, управление зависимостями между работами и визуализация плана на диаграмме Ганта. Клиентское программное обеспечение устанавливать не требуется -- достаточно современного браузера.

Задачами разработки являются:
\begin{itemize}
    \item реализация подсистемы регистрации, авторизации и управления учетными записями пользователей;
    \item разработка модуля создания проектов и управления их основными параметрами;
    \item реализация системы ролей участников проекта: владелец, редактор, наблюдатель;
    \item разработка модуля задач с назначением дат начала и завершения;
    \item реализация механизма зависимостей между задачами с поддержкой четырех типов: финиш-старт, старт-старт, финиш-финиш, старт-финиш;
    \item разработка диаграммы Ганта как основного представления календарного плана проекта;
    \item разработка адаптивного пользовательского интерфейса с минималистичным оформлением, корректно отображающегося на устройствах с различными размерами экрана;
    \item реализация персонализации цветовой схемы интерфейса.
\end{itemize}

\subsection{Требования к функциональности}

Перечень функций системы:

\begin{enumerate}
    \item Регистрация и авторизация. Система должна обеспечивать создание учетных записей с заданием имени пользователя, адреса электронной почты и пароля. Пароль должен соответствовать требованиям: длина не менее восьми символов, наличие хотя бы одной заглавной буквы, одной строчной буквы, одной цифры и одного специального символа, использование только допустимых символов. Форма регистрации также включает поле подтверждения пароля. После успешного входа сервер открывает сеанс работы и записывает идентификатор сеанса в защищенную cookie, которая передается браузером с каждым последующим запросом и используется для аутентификации.

    \item Управление профилем. Пользователь должен иметь возможность просматривать и редактировать данные своей учетной записи: имя, фамилию, отчество и адрес электронной почты.

    \item Настройка цветовой темы. Система должна предоставлять возможность выбора одного из четырех режимов темы оформления: системного (следует настройкам операционной системы), светлого, темного и пользовательского. В пользовательском режиме должна быть доступна настройка цветов отдельных элементов интерфейса: фона, заголовков, цвета задач на диаграмме, цвета стрелок зависимостей и прочих элементов.

    \item Реестр проектов. Система должна отображать список проектов, в которых участвует пользователь, и обеспечивать переход к любому из них.

    \item Создание, редактирование и удаление проекта. Пользователь должен иметь возможность создавать проект с указанием названия, описания и дат начала и завершения. При создании проекта пользователь автоматически получает роль владельца. Редактирование и удаление проекта доступны только владельцу.

    \item Управление участниками. Владелец проекта должен иметь возможность добавлять участников по имени пользователя или ФИО, назначать им роли редактора или наблюдателя, изменять роль и удалять участников из проекта.

    \item Выход из проекта. Любой участник проекта, кроме владельца, должен иметь возможность самостоятельно покинуть проект. После выхода доступ к данным проекта утрачивается.

    \item Создание, редактирование и удаление задач. Пользователи с ролью владельца или редактора должны иметь возможность создавать задачи в рамках проекта с указанием названия, описания, даты начала и даты завершения, редактировать и удалять их. Каждой задаче автоматически присваивается порядковый номер в рамках проекта. При нажатии на полосу задачи на диаграмме Ганта открывается детальная карточка задачи, в которой отображаются и редактируются все поля, а также список зависимостей задачи. Дата начала и дата завершения в карточке ограничены снизу и сверху датами проекта и ограничениями, вытекающими из установленных зависимостей. Порядок задач на диаграмме должен быть изменяем перетаскиванием строк по вертикали.

    \item Управление зависимостями задач. Владелец и редактор должны иметь возможность устанавливать зависимости между задачами. Поддерживаются четыре типа зависимостей: финиш-старт (задача-преемник начинается после завершения предшественника), старт-старт (одновременное начало), финиш-финиш (одновременное завершение) и старт-финиш (редко применяемый тип, при котором завершение преемника привязано к началу предшественника). Система должна исключать создание зависимости задачи от самой себя и дублирующих пар. На диаграмме Ганта должен быть предусмотрен режим добавления связей: в этом режиме на торцах каждой полосы задачи появляются точки соединения, и зависимость создается перетаскиванием от одной точки к другой с автоматическим определением типа по выбранным торцам. Зависимости также должны быть доступны для добавления и удаления из карточки задачи.

    \item Диаграмма Ганта. Система должна отображать задачи проекта в виде диаграммы Ганта: каждая задача представлена горизонтальным отрезком, позиция и длина которого определяются датами начала и завершения. На полосе задачи отображаются ее название и длительность в днях. Зависимости между задачами должны отображаться стрелками соответствующего типа. Даты задачи должны изменяться перетаскиванием непосредственно на диаграмме: горизонтальное перемещение полосы смещает обе даты с сохранением длительности, а перетаскивание левого или правого края изменяет дату начала или дату завершения соответственно. При перемещении или изменении размеров полосы задачи система должна автоматически ограничивать допустимый диапазон дат в соответствии с установленными зависимостями: например, при зависимости финиш-старт задача-преемник не может быть сдвинута раньше даты завершения предшественника. Временная шкала охватывает исключительно период от даты начала до даты завершения проекта. Пользователь должен иметь возможность изменять масштаб отображения в диапазоне от 50 до 400~\% с помощью кнопок панели инструментов, прокрутки колесом мыши с удержанием Ctrl или жестом сведения пальцев на сенсорном экране. Выбранный масштаб должен сохраняться отдельно для каждого проекта. Должна быть предусмотрена возможность блокировки диаграммы: в заблокированном состоянии перетаскивание полос и изменение их размеров недоступны, однако просмотр и открытие карточек задач сохраняется.

    \item Разграничение прав доступа. Предусмотрены три уровня доступа к данным проекта. Владелец -- полный доступ, включая управление составом участников и параметрами проекта. Редактор -- работа с задачами и зависимостями без права изменять состав участников. Наблюдатель -- только чтение данных проекта.

    \item Валидация данных. Система должна выполнять проверку корректности данных, вводимых при регистрации, редактировании профиля, создании и изменении проекта, а также при работе с задачами.
\end{enumerate}

\subsection{Требования к интерфейсу}

Интерфейс реализован как одностраничное веб-приложение (SPA). Страницы переключаются без перезагрузки. После авторизации на всех экранах остаются шапка и боковая панель навигации.

В шапке -- выпадающий список проектов пользователя с постраничной подгрузкой, кнопка создания нового проекта и кнопка добавления задачи в текущий проект. Последняя отображается только пользователям с ролью редактора или владельца.

Боковая панель содержит четыре кнопки перехода между разделами: диаграмма Ганта, профиль, проект, настройки. Панель можно свернуть до иконок, чтобы освободить место для рабочей области.

Основные разделы интерфейса:
\begin{itemize}
    \item страница входа и регистрации -- видна неавторизованным пользователям, после входа становится недоступной;
    \item страница диаграммы Ганта -- основной рабочий экран с интерактивным отображением задач и зависимостей;
    \item страница проекта -- редактирование параметров текущего проекта (название, описание, даты) и управление составом участников;
    \item страница профиля пользователя -- просмотр и редактирование персональных данных учетной записи;
    \item страница настроек -- управление цветовой темой интерфейса.
\end{itemize}

Интерфейс адаптирован под разные размеры экрана -- от смартфонов до ПК. Раскладка перестраивается по ширине. Горизонтальная прокрутка и потеря функциональности на узких экранах исключены.

Оформление минималистичное. Декоративных элементов нет, функциональные зоны разделены, поведение элементов предсказуемо. Назначение каждого элемента понятно без подсказок. Навигация не уходит глубже двух уровней.

Интерактивные элементы откликаются на действия пользователя: видны состояния наведения, нажатия и недоступности. Цвета вынесены в CSS-переменные, поэтому цветовую схему можно менять без правок исходного кода. Шрифты, отступы и элементы управления одинаковы на всех страницах.

Внешний вид страницы входа и регистрации показан на рисунке~\ref{fig:tz_welcome}.

\begin{figure}[H]
\centering

\includegraphics[width=\textwidth]{sc_welcome}
\caption{Страница входа в систему}
\label{fig:tz_welcome}
\end{figure}

При вводе пароля форма регистрации показывает список требований к нему в реальном времени: каждый выполненный пункт подсвечивается сразу. На рисунке~\ref{fig:tz_password_hint} показан вид подсказки при частично введенном пароле.

\begin{figure}[H]
\centering
\includegraphics[width=\textwidth]{sc_password_hint}
\caption{Интерактивная подсказка к паролю при регистрации}
\label{fig:tz_password_hint}
\end{figure}

Рисунок~\ref{fig:tz_gantt} иллюстрирует страницу диаграммы Ганта с отображением задач, зависимостей, линии текущего времени и штриховки истекшей части полос.

\begin{figure}[H]
\centering
\includegraphics[width=\textwidth]{sc_gantt}
\caption{Страница с диаграммой Ганта проекта}
\label{fig:tz_gantt}
\end{figure}

На диаграмме Ганта выводится вертикальная линия текущей даты -- только если она попадает в период проекта. Истекшая часть каждой полосы задачи выделяется штриховкой.

При нажатии на полосу задачи открывается карточка с детальной информацией и элементами управления. На рисунке~\ref{fig:tz_task_card} показан вид карточки задачи.

\begin{figure}[H]
\centering
\includegraphics[width=\textwidth]{sc_task_card}
\caption{Карточка задачи}
\label{fig:tz_task_card}
\end{figure}

Страница проекта с параметрами и составом участников приведена на рисунке~\ref{fig:tz_project_page}.

\begin{figure}[H]
\centering
\includegraphics[width=\textwidth]{sc_project_page}
\caption{Страница проекта}
\label{fig:tz_project_page}
\end{figure}

Страница настроек с элементами управления цветовой темой показана на рисунке~\ref{fig:tz_settings_page}.

\begin{figure}[H]
\centering
\includegraphics[width=\textwidth]{sc_settings_page}
\caption{Страница настроек}
\label{fig:tz_settings_page}
\end{figure}

\subsection{Моделирование вариантов использования}

Модель прецедентов построена по ролям участников \cite{jacobson_rup}. Диаграммы выполнены на языке UML \cite{omg_uml}.

В системе четыре роли. Гость -- незарегистрированный посетитель, ему доступны только формы входа и регистрации. Наблюдатель сверх этого работает с реестром проектов, профилем, настройками темы и просматривает данные тех проектов, в которых он состоит. Редактор наследует права наблюдателя и управляет задачами -- создает, редактирует, удаляет, меняет порядок и зависимости. Владелец обладает всеми правами редактора, а сверх того редактирует и удаляет проект, управляет составом команды и ролями участников.

На основании анализа предметной области в системе реализованы следующие прецеденты:

\begin{enumerate}
    \item Регистрация пользователя.
    \item Вход в систему.
    \item Выход из системы.
    \item Просмотр профиля пользователя.
    \item Редактирование профиля пользователя.
    \item Настройка цветовой темы интерфейса.
    \item Просмотр списка проектов.
    \item Создание проекта.
    \item Редактирование параметров проекта.
    \item Удаление проекта.
    \item Добавление участника в проект.
    \item Изменение роли участника.
    \item Удаление участника из проекта.
    \item Выход из проекта.
    \item Просмотр диаграммы Ганта.
    \item Создание задачи.
    \item Редактирование задачи.
    \item Удаление задачи.
    \item Изменение порядка задач.
    \item Создание зависимости между задачами.
    \item Удаление зависимости между задачами.
\end{enumerate}

Диаграмма прецедентов гостя приведена на рисунке~\ref{fig:uc_guest}.

\begin{figure}[H]
\centering
\includegraphics[width=0.7\textwidth]{use_case_guest}
\caption{Диаграмма прецедентов гостя}
\label{fig:uc_guest}
\end{figure}

На рисунке~\ref{fig:uc_user} показана диаграмма прецедентов наблюдателя проекта.

\begin{figure}[H]
\centering
\includegraphics[width=\textwidth]{use_case_observer}
\caption{Диаграмма прецедентов наблюдателя проекта}
\label{fig:uc_user}
\end{figure}

Рисунок~\ref{fig:uc_editor} содержит диаграмму прецедентов редактора проекта.

\begin{figure}[H]
\centering
\includegraphics[width=\textwidth]{use_case_editor}
\caption{Диаграмма прецедентов редактора проекта}
\label{fig:uc_editor}
\end{figure}

Диаграмма прецедентов владельца проекта приведена на рисунке~\ref{fig:uc_owner}.

\begin{figure}[H]
\centering
\includegraphics[width=\textwidth]{use_case_owner}
\caption{Диаграмма прецедентов владельца проекта}
\label{fig:uc_owner}
\end{figure}

\subsubsection{Прецедент «Регистрация пользователя»}

\textbf{Заинтересованные лица и их требования:} незарегистрированный пользователь хочет создать учетную запись для работы с системой.

\textbf{Предусловие:} пользователь не аутентифицирован в системе.

\textbf{Постусловие:} в системе создана новая учетная запись, пользователь выполнил вход.

\textbf{Основной успешный сценарий:}
\begin{enumerate}
    \item Пользователь нажимает кнопку «Зарегистрироваться» на стартовой странице.
    \item Система открывает модальное окно регистрации с полями: логин, адрес электронной почты, пароль, подтверждение пароля.
    \item Пользователь заполняет поля. При вводе пароля система в реальном времени отображает интерактивный список требований с индикацией выполнения каждого пункта.
    \item Пользователь нажимает кнопку «Зарегистрироваться» в модальном окне.
    \item Система проверяет уникальность логина и адреса электронной почты, соответствие пароля требованиям, совпадение паролей.
    \item Система создает учетную запись с хешированным паролем.
    \item Система выполняет вход и перенаправляет пользователя на страницу диаграммы Ганта.
\end{enumerate}

\textbf{Альтернативный сценарий:}
\begin{enumerate}
    \item На шаге 5 система обнаруживает, что логин или адрес электронной почты уже заняты либо пароль не удовлетворяет требованиям.
    \item Система отображает сообщение об ошибке рядом с соответствующим полем.
    \item Пользователь исправляет данные и повторяет отправку.
\end{enumerate}

\subsubsection{Прецедент «Вход в систему»}

\textbf{Заинтересованные лица и их требования:} зарегистрированный пользователь намерен войти и продолжить работу с системой.

\textbf{Предусловие:} пользователь не аутентифицирован, открыта стартовая страница.

\textbf{Постусловие:} пользователь аутентифицирован, открыт сеанс работы, пользователь перенаправлен на страницу диаграммы Ганта.

\textbf{Основной успешный сценарий:}
\begin{enumerate}
    \item Пользователь нажимает кнопку «Войти» на стартовой странице.
    \item Система открывает модальное окно входа с полями: логин и пароль.
    \item Пользователь вводит данные и нажимает кнопку «Войти».
    \item Система проверяет существование учетной записи и корректность пароля.
    \item Система открывает сеанс работы.
    \item Пользователь перенаправляется на страницу диаграммы Ганта.
\end{enumerate}

\textbf{Альтернативный сценарий:}
\begin{enumerate}
    \item На шаге 4 система обнаруживает, что учетная запись не найдена или пароль неверен.
    \item Система отображает сообщение об ошибке.
    \item Пользователь остается на стартовой странице и может повторить попытку.
\end{enumerate}

\subsubsection{Прецедент «Выход из системы»}

\textbf{Заинтересованные лица и их требования:} пользователь намерен завершить текущий сеанс.

\textbf{Предусловие:} пользователь вошел в систему.

\textbf{Постусловие:} сеанс работы завершен, пользователь перенаправлен на стартовую страницу.

\textbf{Основной успешный сценарий:}
\begin{enumerate}
    \item Пользователь нажимает кнопку выхода в интерфейсе приложения.
    \item Система завершает сеанс работы пользователя.
    \item Система перенаправляет пользователя на стартовую страницу.
\end{enumerate}

\subsubsection{Прецедент «Просмотр профиля пользователя»}

\textbf{Заинтересованные лица и их требования:} пользователь хочет посмотреть данные своей учетной записи.

\textbf{Предусловие:} пользователь авторизован.

\textbf{Постусловие:} на экране отображены данные профиля пользователя.

\textbf{Основной успешный сценарий:}
\begin{enumerate}
    \item Пользователь выбирает раздел «Профиль» в боковой панели навигации.
    \item Система отображает страницу профиля: имя пользователя, имя, фамилию, отчество и адрес электронной почты.
\end{enumerate}

\subsubsection{Прецедент «Редактирование профиля пользователя»}

\textbf{Заинтересованные лица и их требования:} пользователь собирается обновить данные учетной записи.

\textbf{Предусловие:} пользователь аутентифицирован, открыта страница собственного профиля.

\textbf{Постусловие:} обновленные данные сохранены в базе данных.

\textbf{Основной успешный сценарий:}
\begin{enumerate}
    \item Пользователь переходит на страницу своего профиля.
    \item Система отображает форму редактирования с полями: имя, фамилия, отчество, адрес электронной почты. Поле имени пользователя недоступно для изменения.
    \item Пользователь изменяет нужные поля и подтверждает сохранение.
    \item Система проверяет корректность данных.
    \item Система сохраняет изменения и отображает обновленные данные.
\end{enumerate}

\textbf{Альтернативный сценарий:}
\begin{enumerate}
    \item На шаге 4 система обнаруживает, что адрес электронной почты уже используется другим пользователем.
    \item Система отображает сообщение об ошибке.
    \item Пользователь корректирует данные и повторяет сохранение.
\end{enumerate}

\subsubsection{Прецедент «Настройка цветовой темы интерфейса»}

\textbf{Заинтересованные лица и их требования:} пользователь хочет сменить цветовую тему интерфейса.

\textbf{Предусловие:} открыт авторизованный сеанс.

\textbf{Постусловие:} выбранная тема применена к интерфейсу.

\textbf{Основной успешный сценарий:}
\begin{enumerate}
    \item Пользователь переходит в раздел «Настройки» через боковую панель навигации.
    \item Система отображает страницу выбора режима темы: системный, светлый, темный, пользовательский.
    \item Пользователь выбирает нужный режим. В пользовательском режиме дополнительно отображаются элементы выбора цвета для отдельных компонентов интерфейса.
    \item Система применяет тему немедленно без перезагрузки страницы.
\end{enumerate}

\subsubsection{Прецедент «Просмотр списка проектов»}

\textbf{Заинтересованные лица и их требования:} пользователю нужно перейти к нужному проекту.

\textbf{Предусловие:} пользователь вошел в систему.

\textbf{Постусловие:} пользователю отображен список проектов, в которых он является участником.

\textbf{Основной успешный сценарий:}
\begin{enumerate}
    \item Пользователь нажимает на выпадающий список проектов в шапке приложения.
    \item Система загружает и отображает список проектов пользователя с постраничной подгрузкой.
    \item Пользователь выбирает проект для перехода к нему.
\end{enumerate}

\subsubsection{Прецедент «Создание проекта»}

\textbf{Заинтересованные лица и их требования:} пользователь хочет создать новый проект.

\textbf{Предусловие:} пользователь авторизован.

\textbf{Постусловие:} в системе создан новый проект, пользователь добавлен в него с ролью владельца.

\textbf{Основной успешный сценарий:}
\begin{enumerate}
    \item Пользователь нажимает кнопку создания проекта в шапке приложения.
    \item Система отображает форму с полями: название, описание, дата начала, дата завершения.
    \item Пользователь заполняет поля и подтверждает создание.
    \item Система проверяет корректность данных.
    \item Система создает проект и назначает пользователю роль владельца.
    \item Система переключает контекст на созданный проект и переходит к его диаграмме Ганта.
\end{enumerate}

\textbf{Альтернативный сценарий:}
\begin{enumerate}
    \item На шаге 4 система обнаруживает незаполненное обязательное поле или некорректные даты (дата завершения раньше даты начала).
    \item Система отображает сообщение об ошибке.
    \item Пользователь исправляет данные и повторяет создание.
\end{enumerate}

\subsubsection{Прецедент «Редактирование параметров проекта»}

\textbf{Заинтересованные лица и их требования:} владелец собирается изменить основные параметры проекта.

\textbf{Предусловие:} пользователь аутентифицирован, пользователь является владельцем проекта, открыта страница проекта.

\textbf{Постусловие:} обновленные параметры проекта сохранены в базе данных.

\textbf{Основной успешный сценарий:}
\begin{enumerate}
    \item Владелец переходит на страницу проекта через боковую панель навигации.
    \item Система отображает форму с текущими значениями: название, описание, дата начала, дата завершения.
    \item Владелец изменяет нужные поля и сохраняет изменения.
    \item Система проверяет корректность данных.
    \item Система сохраняет изменения.
\end{enumerate}

\textbf{Альтернативный сценарий:}
\begin{enumerate}
    \item На шаге 4 система обнаруживает незаполненное обязательное поле или некорректный порядок дат.
    \item Система отображает сообщение об ошибке.
    \item Владелец исправляет данные и повторяет сохранение.
\end{enumerate}

\subsubsection{Прецедент «Удаление проекта»}

\textbf{Заинтересованные лица и их требования:} владелец намерен удалить проект со всеми данными.

\textbf{Предусловие:} пользователь аутентифицирован, пользователь является владельцем проекта.

\textbf{Постусловие:} проект со всеми задачами, зависимостями и записями об участниках удален из системы.

\textbf{Основной успешный сценарий:}
\begin{enumerate}
    \item Владелец на странице проекта нажимает кнопку удаления проекта.
    \item Система отображает модальное окно подтверждения с пояснением о необратимости операции.
    \item Владелец подтверждает удаление.
    \item Система удаляет проект и все связанные с ним данные.
    \item Если у пользователя остаются другие проекты, система переключает контекст на один из них; в противном случае создается новый пустой проект, и пользователь перенаправляется на его диаграмму Ганта.
\end{enumerate}

\textbf{Альтернативный сценарий:}
\begin{enumerate}
    \item На шаге 3 владелец закрывает модальное окно без подтверждения.
    \item Система отменяет операцию, проект остается без изменений.
\end{enumerate}

\subsubsection{Прецедент «Добавление участника в проект»}

\textbf{Заинтересованные лица и их требования:} владелец хочет добавить нового участника в проект.

\textbf{Предусловие:} пользователь аутентифицирован, пользователь является владельцем проекта, открыта страница проекта.

\textbf{Постусловие:} новый участник добавлен в проект с указанной ролью.

\textbf{Основной успешный сценарий:}
\begin{enumerate}
    \item Владелец открывает раздел управления участниками на странице проекта.
    \item Система отображает текущий состав участников и форму добавления.
    \item Владелец вводит имя пользователя или ФИО добавляемого пользователя в поле поиска.
    \item Система отображает список совпадений по имени пользователя и полному имени.
    \item Владелец выбирает пользователя из списка и назначает ему роль: редактор или наблюдатель.
    \item Система добавляет пользователя в проект с выбранной ролью.
\end{enumerate}

\textbf{Альтернативный сценарий:}
\begin{enumerate}
    \item На шаге 4 система не находит пользователей по введенному запросу.
    \item Система отображает пустой список результатов.
    \item Владелец уточняет запрос и повторяет поиск.
\end{enumerate}

\subsubsection{Прецедент «Изменение роли участника»}

\textbf{Заинтересованные лица и их требования:} владелец намерен поменять роль одному из участников.

\textbf{Предусловие:} пользователь аутентифицирован, пользователь является владельцем проекта, открыта страница проекта.

\textbf{Постусловие:} роль участника обновлена в базе данных.

\textbf{Основной успешный сценарий:}
\begin{enumerate}
    \item Владелец открывает раздел управления участниками на странице проекта.
    \item Система отображает список участников с текущими ролями.
    \item Владелец выбирает новую роль для нужного участника.
    \item Система сохраняет изменение.
\end{enumerate}

\subsubsection{Прецедент «Удаление участника из проекта»}

\textbf{Заинтересованные лица и их требования:} владелец собирается исключить участника из проекта.

\textbf{Предусловие:} пользователь аутентифицирован, пользователь является владельцем проекта, открыта страница проекта.

\textbf{Постусловие:} участник удален из проекта, доступ к его данным утрачен.

\textbf{Основной успешный сценарий:}
\begin{enumerate}
    \item Владелец открывает раздел управления участниками на странице проекта.
    \item Система отображает список участников.
    \item Владелец нажимает кнопку удаления рядом с нужным участником.
    \item Система отображает модальное окно подтверждения.
    \item Владелец подтверждает удаление.
    \item Система удаляет запись об участии пользователя в проекте.
\end{enumerate}

\subsubsection{Прецедент «Выход из проекта»}

\textbf{Заинтересованные лица и их требования:} участник с ролью редактора или наблюдателя хочет самостоятельно покинуть проект.

\textbf{Предусловие:} пользователь аутентифицирован, пользователь является участником проекта с ролью редактора или наблюдателя.

\textbf{Постусловие:} запись об участии удалена, доступ к данным проекта утрачен.

\textbf{Основной успешный сценарий:}
\begin{enumerate}
    \item Участник открывает раздел управления участниками на странице проекта.
    \item Система отображает кнопку «Покинуть проект».
    \item Участник нажимает кнопку.
    \item Система удаляет запись об участии и перенаправляет пользователя.
\end{enumerate}

\subsubsection{Прецедент «Просмотр диаграммы Ганта»}

\textbf{Заинтересованные лица и их требования:} участник намерен просмотреть план работ на диаграмме Ганта.

\textbf{Предусловие:} пользователь аутентифицирован и является участником хотя бы одного проекта.

\textbf{Постусловие:} на экране отображена интерактивная диаграмма Ганта с задачами, зависимостями и временной шкалой.

\textbf{Основной успешный сценарий:}
\begin{enumerate}
    \item Пользователь выбирает проект в выпадающем списке шапки и переходит к разделу диаграммы через боковую панель.
    \item Система загружает задачи и зависимости проекта.
    \item Система строит временную шкалу в пределах дат проекта и отображает каждую задачу горизонтальной полосой. Позиция и длина полосы определяются датами начала и завершения задачи, на полосе выводятся название задачи и ее длительность в днях.
    \item Система отображает зависимости стрелками, вертикальную линию текущего времени (если текущая дата входит в период проекта) и штриховку истекшей части каждой полосы.
    \item Пользователь может изменять масштаб диаграммы от 50 до 400~\% с помощью кнопок панели инструментов, прокрутки колесом мыши с удержанием Ctrl или жестом сведения пальцев. Масштаб сохраняется отдельно для каждого проекта.
\end{enumerate}

\subsubsection{Прецедент «Создание задачи»}

\textbf{Заинтересованные лица и их требования:} редактор или владелец собирается добавить задачу в план проекта.

\textbf{Предусловие:} пользователь аутентифицирован, пользователь имеет роль редактора или владельца в текущем проекте.

\textbf{Постусловие:} задача создана, ей присвоен порядковый номер, она отображается на диаграмме Ганта.

\textbf{Основной успешный сценарий:}
\begin{enumerate}
    \item Пользователь нажимает кнопку добавления задачи в шапке приложения.
    \item Система отображает форму с полями: название, описание, дата начала, дата завершения.
    \item Пользователь заполняет поля и подтверждает создание.
    \item Система проверяет корректность данных.
    \item Система создает задачу и присваивает ей следующий порядковый номер в рамках проекта.
    \item Диаграмма Ганта обновляется: новая задача появляется в конце списка.
\end{enumerate}

\textbf{Альтернативный сценарий:}
\begin{enumerate}
    \item На шаге 4 система обнаруживает незаполненное обязательное поле или даты, выходящие за пределы дат проекта.
    \item Система отображает сообщение об ошибке.
    \item Пользователь исправляет данные и повторяет создание.
\end{enumerate}

\subsubsection{Прецедент «Редактирование задачи»}

\textbf{Заинтересованные лица и их требования:} редактор или владелец хочет отредактировать параметры задачи.

\textbf{Предусловие:} пользователь аутентифицирован, пользователь имеет роль редактора или владельца, диаграмма не заблокирована, задача существует в проекте.

\textbf{Постусловие:} измененные данные задачи сохранены, диаграмма Ганта обновлена.

\textbf{Основной успешный сценарий (через карточку задачи):}
\begin{enumerate}
    \item Пользователь нажимает на полосу задачи на диаграмме Ганта.
    \item Система открывает карточку задачи с полями: название, описание, дата начала, дата завершения. Поля дат ограничены датами проекта и ограничениями из установленных зависимостей.
    \item Пользователь изменяет нужные поля и сохраняет.
    \item Система проверяет корректность данных.
    \item Система сохраняет изменения, диаграмма Ганта обновляется.
\end{enumerate}

\textbf{Альтернативный сценарий (перетаскивание на диаграмме):}
\begin{enumerate}
    \item Пользователь перетаскивает полосу задачи горизонтально: смещаются обе даты с сохранением длительности. Или перетаскивает левый край полосы для изменения даты начала, правый -- для изменения даты завершения.
    \item Система ограничивает допустимый диапазон перемещения согласно зависимостям.
    \item После отпускания диаграмма отображает новые даты. В панели инструментов становятся доступны кнопки подтверждения и отката изменений.
    \item Пользователь нажимает кнопку подтверждения.
    \item Система отправляет измененные даты на сервер и обновляет диаграмму.
\end{enumerate}

\textbf{Сценарий отката (перетаскивание на диаграмме):}
\begin{enumerate}
    \item На шаге 4 пользователь нажимает кнопку отката.
    \item Система возвращает задачи к исходным датам без обращения к серверу.
\end{enumerate}

\subsubsection{Прецедент «Удаление задачи»}

\textbf{Заинтересованные лица и их требования:} редактор или владелец намерен удалить задачу из проекта.

\textbf{Предусловие:} пользователь аутентифицирован, пользователь имеет роль редактора или владельца, задача существует в проекте.

\textbf{Постусловие:} задача удалена из проекта вместе со всеми ее зависимостями, диаграмма Ганта обновлена.

\textbf{Основной успешный сценарий:}
\begin{enumerate}
    \item Пользователь открывает карточку задачи нажатием на полосу диаграммы.
    \item Система отображает карточку с кнопкой удаления задачи.
    \item Пользователь нажимает кнопку удаления.
    \item Система отображает модальное окно подтверждения.
    \item Пользователь подтверждает удаление.
    \item Система удаляет задачу и все связанные зависимости.
    \item Диаграмма Ганта обновляется.
\end{enumerate}

\subsubsection{Прецедент «Изменение порядка задач»}

\textbf{Заинтересованные лица и их требования:} редактор или владелец хочет переставить задачи в другом порядке на диаграмме.

\textbf{Предусловие:} пользователь аутентифицирован, пользователь имеет роль редактора или владельца, в проекте существует не менее двух задач, диаграмма не заблокирована.

\textbf{Постусловие:} пользователь подтвердил изменения, порядковые номера задач обновлены на сервере, диаграмма отображает задачи в новом порядке.

\textbf{Основной успешный сценарий:}
\begin{enumerate}
    \item Пользователь захватывает строку задачи за ручку перетаскивания в левой части строки.
    \item Пользователь перемещает строку на новую позицию.
    \item После отпускания диаграмма отображает обновленный порядок. В панели инструментов становятся доступны кнопки подтверждения и отката изменений.
    \item Пользователь нажимает кнопку подтверждения.
    \item Система отправляет измененный порядок задач на сервер и обновляет порядковые номера.
\end{enumerate}

\textbf{Сценарий отката:}
\begin{enumerate}
    \item На шаге 4 пользователь нажимает кнопку отката.
    \item Система возвращает задачи к исходному порядку без обращения к серверу.
\end{enumerate}

\subsubsection{Прецедент «Создание зависимости между задачами»}

\textbf{Заинтересованные лица и их требования:} редактор или владелец собирается создать зависимость между двумя задачами.

\textbf{Предусловие:} пользователь аутентифицирован, пользователь имеет роль редактора или владельца, в проекте существует не менее двух задач.

\textbf{Постусловие:} зависимость сохранена, на диаграмме Ганта появилась соответствующая стрелка.

\textbf{Основной успешный сценарий (через режим связей на диаграмме):}
\begin{enumerate}
    \item Пользователь активирует режим добавления связей на панели инструментов диаграммы.
    \item Система переходит в режим связей: на левом и правом торцах каждой полосы задачи появляются точки соединения.
    \item Пользователь нажимает на точку соединения задачи-предшественника и тянет до точки задачи-преемника.
    \item Система определяет тип зависимости по выбранным торцам: левый торец -- старт, правый -- финиш. Сочетание торцов однозначно задает тип связи (финиш-старт, старт-старт, финиш-финиш, старт-финиш).
    \item Система сохраняет зависимость и отображает стрелку на диаграмме.
\end{enumerate}

\textbf{Альтернативный сценарий (через карточку задачи):}
\begin{enumerate}
    \item Пользователь открывает карточку задачи.
    \item В разделе зависимостей пользователь выбирает связанную задачу из списка и задает тип зависимости.
    \item Система сохраняет зависимость и обновляет диаграмму.
\end{enumerate}

\textbf{Сценарий ошибки:}
\begin{enumerate}
    \item Пользователь пытается создать зависимость задачи от самой себя или дублирующую пару с тем же типом.
    \item Система отклоняет операцию и отображает сообщение об ошибке.
\end{enumerate}

\subsubsection{Прецедент «Удаление зависимости между задачами»}

\textbf{Заинтересованные лица и их требования:} редактор или владелец хочет удалить зависимость между задачами.

\textbf{Предусловие:} пользователь аутентифицирован, пользователь имеет роль редактора или владельца, зависимость между задачами существует.

\textbf{Постусловие:} зависимость удалена, стрелка на диаграмме Ганта исчезла.

\textbf{Основной успешный сценарий:}
\begin{enumerate}
    \item Пользователь открывает карточку задачи.
    \item В разделе зависимостей пользователь нажимает кнопку удаления рядом с нужной зависимостью.
    \item Система удаляет зависимость и обновляет диаграмму.
\end{enumerate}

\subsection{Требования к надежности}

Целостность данных обеспечивается тем, что все данные хранятся в MySQL, а операции, затрагивающие несколько связанных таблиц (например, создание проекта с одновременным добавлением владельца в таблицу участников), выполняются в одной транзакции. Появление частично сохраненных или противоречивых записей исключено.

Сбой при выполнении операции не повреждает уже сохраненные данные -- состояние базы остается прежним, а клиент получает сообщение об ошибке с причиной. Граничные случаи на стороне пользователя -- повторная отправка формы, попытка добавить уже состоящего в проекте участника, даты в неверном порядке -- обрабатываются без нарушения работоспособности.

Сервер проверяет все входящие данные. К проверкам относятся формат дат, длина строковых полей, уникальность пар (проект, пользователь) в таблице участников, отсутствие самозависимостей и дублирующих пар в таблице зависимостей.

Резервные копии базы создаются штатными средствами MySQL, восстановление выполняется из них без дополнительного инструментария.

\subsection{Требования к информационной безопасности}

Все программные интерфейсы, кроме регистрации и входа, требуют действующего сеанса. Запросы без сеанса отклоняются с кодом 401. Для изменяющих операций требуется еще и действующий CSRF-токен -- сервер отдает его клиенту отдельной cookie.

Пароли хранятся в виде хеша, алгоритм выбирается из встроенных механизмов Django. При входе сравниваются хеши, открытый пароль нигде не сохраняется.

Разграничение прав доступа выполнено на основе ролевой модели. Права проверяются на сервере до обработки каждого запроса, который требует повышенных привилегий. К данным проектов, в которых пользователь не состоит, доступа нет вообще -- ни на чтение, ни на изменение. Каждый пользователь работает только со своей учетной записью, а возможность отредактировать чужой профиль, задачу или параметры проекта на сервере исключена.

Для защиты от внедрения SQL-кода \cite{owasp_top10, art_web_security} все запросы к базе проходят через ORM Django с параметризованными запросами. Конкатенация пользовательского ввода в SQL-выражения не применяется.

При выходе сеанс закрывается на сервере, а cookie с идентификатором удаляется. Попытка использовать закрытый сеанс отклоняется.

\subsection{Требования к программной и технической совместимости}

Операционная система. Серверная часть работает под Linux или Windows при наличии Python 3.10+ и MySQL 8.0+. Клиентская часть запускается в браузере -- поддерживаются Windows, macOS, Linux, Android и iOS.

Аппаратные требования к серверу:
\begin{itemize}
    \item процессор с тактовой частотой не ниже 2~ГГц;
    \item оперативная память объемом не менее 2~ГБ;
    \item наличие свободного дискового пространства для размещения программных файлов и базы данных;
    \item постоянный доступ к сети.
\end{itemize}

Программное обеспечение. Стек серверной части:
\begin{itemize}
    \item язык программирования Python 3.10+;
    \item веб-фреймворк Django 6.0+;
    \item библиотека Django REST Framework 3.16+ для построения программного интерфейса;
    \item СУБД MySQL 8.0+.
\end{itemize}

На стороне клиента используются:
\begin{itemize}
    \item библиотеки React 19+;
    \item языки HTML5, CSS3 и JavaScript.
\end{itemize}

Для работы с клиентской частью пользователю достаточно современного браузера с поддержкой HTML5 и JavaScript \cite{mdn_docs}: Google Chrome, Mozilla Firefox, Microsoft Edge, Safari или иной совместимый.

Сетевые требования. Клиентская и серверная части обмениваются данными в формате JSON по протоколу HTTP. В производственной среде канал связи должен быть защищен HTTPS.

\subsection{Требования к оформлению документации}

Разработка программной документации и программного изделия должна производиться согласно ГОСТ 19.102--77 \cite{gost19102} и ГОСТ 34.601--90 \cite{gost34601}. Качество разрабатываемого программного обеспечения оценивается по модели характеристик качества, установленной ГОСТ Р ИСО/МЭК 25010--2015 \cite{gost25010}.
